\documentclass{article}
\usepackage[utf8]{inputenc}
\usepackage{amsmath}
\usepackage{amsfonts}
\usepackage{amssymb}
\usepackage{geometry}
\usepackage{tcolorbox}
\usepackage{hyperref}
\geometry{a4paper, margin=1in}

\title{Module 08: Production Pipelines \& Collaborative Git Workflows - Part 4}
\author{Cybernauts 2026 Datathon Campaign}
\date{May 2026}

\begin{document}

\maketitle

\begin{abstract}
The difference between a podium finish and an unranked exit in a fast-paced datathon often comes down to how a team spends the first hour of the competition. Teams that waste time re-writing standard cross-validation splits, model configurations, and file exporters fall behind immediately. This note details the architectural design of a unified, plug-and-play "Secret Weapon" pipeline template, engineered to ingest raw data and export optimized predictions within 20 minutes of kickoff.
\end{abstract}

\tableofcontents
\newpage

\section{The Philosophy of Zero Boilerplate}
When the competitive window opens and the data files are released, your team is fighting a battle against the clock. Writing standard orchestration loops, setting up evaluation matrices, and debugging shape mismatches under high stress is an unnecessary waste of energy.

\subsection{The 20-Minute Leaderboard Target}
Your operational goal is clear: within 20 minutes of the datasets being made available, your team must lock in its very first functional baseline submission on the leaderboard. Achieving this rapid turnaround requires consolidating a week's worth of engineering, validation, and optimization code into a single, fully integrated, reusable master script template.

With the boilerplate baseline secure, your team can spend the remaining hours of the datathon focusing on feature extraction and advanced model exploration.

\section{Anatomy of the Master Template}
The "Secret Weapon" template must act as a self-contained automation loop. To balance speed with elite performance, the master architecture must integrate five core components into a single file:

\begin{itemize}
    \item \textbf{A Standard Configuration Boundary:} A localized control panel at the top of the file to manage data paths, seed parameters, and execution flags.
    \item \textbf{A Stratified 5-Fold Cross-Validation Loop:} A cross-validation engine designed to track out-of-fold predictions, keeping validation metrics highly accurate and stable.
    \item \textbf{An Integrated Optuna Search Engine:} A Bayesian parameter exploration engine that can be toggled on or off to tune model hyperparameters automatically.
    \item \textbf{Multi-Model LightGBM and XGBoost Templates:} Parallel gradient boosting structures pre-configured with early-stopping callbacks to prevent overfitting.
    \item \textbf{An Automated Post-Processing Suite:} A pipeline step that automatically handles probability blending and threshold scanning to maximize classification metrics.
\end{itemize}

\section{The Blueprint Implementation}
The layout below provides the full, ready-to-deploy Python template framework for binary classification tasks. It serves as your team's core codebase asset for immediate deployment.

\begin{verbatim}
import os
import numpy as np
import pandas as pd
import optuna
from lightgbm import LGBMClassifier
from xgboost import XGBClassifier
from sklearn.model_selection import StratifiedKFold
from sklearn.metrics import log_loss, f1_score

# ==========================================
# 1. CORE PIPELINE CONFIGURATION CONTROL
# ==========================================
class Config:
    TRAIN_PATH = "data/train.csv"
    TEST_PATH = "data/test.csv"
    TARGET_COL = "target"
    SEED = 2026
    FOLDS = 5
    RUN_TUNING = False  # Toggle for Optuna optimization

# ==========================================
# 2. SEAMLESS PIPELINE EXECUTION ENGINE
# ==========================================
def run_pipeline():
    print("[*] Ingesting Datasets...")
    train = pd.read_csv(Config.TRAIN_PATH)
    test = pd.read_csv(Config.TEST_PATH)

    # Identify feature blocks dynamically
    features = [c for c in train.columns if c not in [Config.TARGET_COL, 'id']]
    X = train[features]
    y = train[Config.TARGET_COL]
    X_test = test[features]

    # Initialize multi-fold prediction structures
    oof_lgb = np.zeros(len(train))
    oof_xgb = np.zeros(len(train))
    test_lgb = np.zeros(len(test))
    test_xgb = np.zeros(len(test))

    skf = StratifiedKFold(n_splits=Config.FOLDS, shuffle=True, random_state=Config.SEED)

    # ==========================================
    # 3. INTERTWINED CROSS-VALIDATION LOOP
    # ==========================================
    for fold, (train_idx, val_idx) in enumerate(skf.split(X, y)):
        print(f"--- Training Fold {fold+1}/{Config.FOLDS} ---")
        X_tr, y_tr = X.iloc[train_idx], y.iloc[train_idx]
        X_va, y_va = X.iloc[val_idx], y.iloc[val_idx]

        # Branch A: LightGBM Core Execution
        model_lgb = LGBMClassifier(n_estimators=1000, random_state=Config.SEED, verbosity=-1)
        model_lgb.fit(X_tr, y_tr, eval_set=[(X_va, y_va)], callbacks=[optuna.integration.lightgbm.early_stopping(50, verbose=False)] if hasattr(optuna.integration, 'lightgbm') else [])
        oof_lgb[val_idx] = model_lgb.predict_proba(X_va)[:, 1]
        test_lgb += model_lgb.predict_proba(X_test)[:, 1] / Config.FOLDS

        # Branch B: XGBoost Core Execution
        model_xgb = XGBClassifier(n_estimators=1000, random_state=Config.SEED, early_stopping_rounds=50)
        model_xgb.fit(X_tr, y_tr, eval_set=[(X_va, y_va)], verbose=False)
        oof_xgb[val_idx] = model_xgb.predict_proba(X_va)[:, 1]
        test_xgb += model_xgb.predict_proba(X_test)[:, 1] / Config.FOLDS

    print(f"[*] LGB OOF LogLoss: {log_loss(y, oof_lgb):.5f}")
    print(f"[*] XGB OOF LogLoss: {log_loss(y, oof_xgb):.5f}")

    # ==========================================
    # 4. ENSEMBLING & BLENDING MATRIX
    # ==========================================
    final_oof = 0.6 * oof_lgb + 0.4 * oof_xgb
    final_test = 0.6 * test_lgb + 0.4 * test_xgb
    print(f"[*] Blended OOF LogLoss: {log_loss(y, final_oof):.5f}")

    # ==========================================
    # 5. METRIC CALIBRATION & SCANNING LOOP
    # ==========================================
    best_thresh = 0.5
    best_f1 = 0.0
    for thresh in np.arange(0.01, 1.00, 0.01):
        score = f1_score(y, (final_oof >= thresh).astype(int))
        if score > best_f1:
            best_f1 = score
            best_thresh = thresh

    print(f"[!] Calibration Locked. Best Threshold: {best_thresh:.2f} | F1: {best_f1:.5f}")

    # Export structural submission file
    submission = pd.DataFrame({'id': test.get('id', test.index), 'target': (final_test >= best_thresh).astype(int)})
    submission.to_csv("submission.csv", index=False)
    print("[+] Submission File Successfully Secured.")

if __name__ == "__main__":
    run_pipeline()
\end{verbatim}

\newpage
\section{Practice Problems}

\textbf{P1:} In a classification datathon, the evaluation metric is changed from Macro F1-Score to Area Under the ROC Curve (ROC-AUC). State what adjustments must be made to the master template's post-processing section to accommodate this change.
\\
\textit{Answer: The threshold scanning calibration loop must be completely skipped or removed. The ROC-AUC metric evaluates a model's ability to rank risk probabilities across continuous distributions, rather than hard binary outputs. For the final submission, the continuous blended probability array (\texttt{final\_test}) must be exported directly into the output file, completely bypassing any binary conversion steps.}

\newline
\textbf{P2:} A competitor copies the template but places the probability blending weights calculation ($0.6 \times \text{LGB} + 0.4 \times \text{XGB}$) inside the K-Fold loop, computing the weights fold-by-fold rather than globally across the complete vectors. Explain whether this change introduces an architectural defect or changes the final output predictions.
\\
\textit{Answer: This adjustment does not introduce an architectural defect, and the final mathematical output remains unchanged. Because linear addition and scalar multiplication are distributive operations over vector spaces, computing a weighted linear blend fold-by-fold on localized slices yields the exact same numerical array as computing the operation globally across the concatenated out-of-fold vectors.}

\newline
\textbf{P3:} Your team plugs a raw datathon dataset into the master template. Within 12 minutes, the script finishes execution and outputs a submission file, but the cross-validation logs show an extremely poor F1-score of $0.12$. Explain how your team should systematically triage this situation under hackathon conditions.
\\
\textit{Answer: The master template has done its job perfectly by validating the execution path. Since the code runs without throwing errors, you can rule out compilation bugs or formatting issues. The poor F1-score indicates an issue with the data features. Your team should immediately begin feature troubleshooting steps: check for skewed numeric values that need log transformations, identify high-cardinality strings that require encoding, or look for extreme class imbalances that need custom weighting tweaks inside the gradient boosting model configurations.}

\end{document}
